\usepackage{color}
\usepackage{graphicx, gensymb}
\usepackage{amsfonts, cancel, amssymb}
\usepackage{amsmath, amsthm, array, esint}
\usepackage{bm} % bold math symbols, e.g. \bm{\sigma} etc.
\usepackage{booktabs, multirow}  % schönere Tabellen
%\usepackage{scrlayer-scrpage}    % Manipulation des Seitenstils
%\pagestyle{scrheadings}  % Stil für die Seite setzen
\usepackage{tabularx}
\usepackage{setspace}
%\automark{section}  % Abschnittsnamen als Seitenbeschriftung 
%\setheadsepline{.5pt}    % Kopzeile durch Linie abtrennen
\usepackage[hidelinks]{hyperref}  % Links und weitere PDF-Features
\usepackage{typearea, tikz, pgffor, verbatim}
\usetikzlibrary{calc, arrows.meta,arrows, graphs, quotes, positioning, angles, babel}
\usepackage[cache=false, outputdir=build]{minted}
\usemintedstyle{vs}
\usepackage{placeins} % provides \FloatBarrier

\definecolor{bg}{rgb}{0.9,0.9,0.9}
\newcommand\numberthis{\addtocounter{equation}{1}\tag{\theequation}}
\usepackage{svg}
\usepackage{siunitx}
\usepackage{physics}
\usepackage{tensor}
\usepackage{slashed}
\usepackage{bbm}
\usepackage[compat=1.1.0]{tikz-feynman}

\usepackage{caption}
\captionsetup{
    % width=.8\textwidth,
    % indention=1cm,
    margin=0.5cm,
    format=plain,
    labelfont=bf
}
%\usepackage{subcaption}
%\subcaptionsetup{
%    indention=0cm,
%    format=hang,
%    margin=0.5cm,
%    labelfont=normalfont
%}
\usepackage[english,capitalise]{cleveref}
\usepackage[
    backend=biber,
    sorting=none,
    style=phys,
    giveninits=true,
    sortcites=true,
    citestyle=numeric-comp,
    % url=false,
    % doi=false,
    biblabel=brackets,
    % eprint=false,
    maxcitenames=3]{biblatex}
\usepackage{csquotes}
%\addbibresource{bibliography.bib}

\usepackage{mathtools}
% use \abs for absolute values
% use \abs for \left ... \right version and \abs[\bigg for \bigg version etc.
\let\abs\undefined
\let\bra\undefined
\let\braket\undefined
\DeclarePairedDelimiter\abs{\vert}{\vert}
\DeclarePairedDelimiter\kla{(}{)}
\DeclarePairedDelimiter\klb{[}{]}
\DeclarePairedDelimiter\klc{\{}{\}}
\DeclarePairedDelimiter\bra{\langle}{\rangle}
\DeclarePairedDelimiterX\brb[2]{\langle #1}{\rangle}{\delimsize\vert #2 }
\DeclarePairedDelimiterX\brc[3]{\langle #1}{\rangle}{\delimsize\vert #2 \delimsize\vert #3 }

\renewcommand{\i}{\text{i}}
\newcommand{\e}{\text{e}}
\DeclareMathOperator{\sgn}{sgn}
\newcommand{\kom}[1]{\overset{\substack{#1 \\ |}}{=}}
\newcommand{\koml}[2]{\overset{\substack{#1 \\ |}}{#2}}
\newcommand{\intx}[4]{\int_{#1}^{#2} #3 \,\mathrm{d}#4}
\newcommand{\intr}[2]{\int_{\mathbb{R}} #1 \, \mathrm{d}#2}
\newcommand{\intrnd}[3]{\int_{\mathbb{R}^{#1}} #2 \, \mathrm{d}^{#1}#3}
\newcommand{\intrpi}[2]{\int_{\mathbb{R}} #1 \, \frac{\mathrm{d}#2}{2\pi}}
\newcommand{\intrndpi}[3]{\int_{\mathbb{R}^{#1}} #2 \, \frac{\mathrm{d}^{#1}#3}{(2\pi)^{#1}}}
\newcommand{\oper}[2]{\vert #1 \rangle \langle #2 \vert}
\newcommand{\Text}[1]{\text{#1}}    % this is relevant because \text does not autocomplete to the default '\text{@}' but rather '\textbf' etc.
\newcommand{\powe}[1]{\e^{#1}}
\newcommand{\pow}[2]{{#1}^{#2}}
\newcommand{\Ket}[1]{\vert #1 \rangle}
\newcommand{\Bra}[1]{\langle #1 \vert}
\newcommand*{\Fourier}{\textsc{Fourier}\ }
\newcommand*{\F}{\mathcal{F}}
\DeclareMathOperator{\arsinh}{arsinh}
\DeclareMathOperator{\arcosh}{arcosh}
\DeclareMathOperator{\artanh}{artanh}
\DeclareMathOperator{\sinc}{sinc}
\DeclareMathOperator{\cscc}{cscc}
\DeclareMathOperator{\sinhc}{sinhc}
\DeclareMathOperator{\cschc}{cschc}
\newcommand{\conv}{\mathop{\scalebox{3.5}{\raisebox{-0.38ex}{\makebox[0.5\width][c]{$\ast$}}}}\limits}
\newcommand*{\unity}{\text{\usefont{U}{bbold}{m}{n}1}}   % operator one from :: https://tex.stackexchange.com/questions/566780/import-mathbb1-character-from-the-unicode-math-package

\renewcommand{\d}{\text{d}}
\renewcommand{\r}{\text{r}}
\renewcommand{\t}{\text{t}}
\renewcommand{\a}{\text{a}}
\renewcommand{\o}{\text{o}}
\renewcommand{\F}{\text{F}}
\renewcommand{\c}{\text{c}}
\newcommand{\A}{\text{A}}
\newcommand{\B}{\text{B}}
\newcommand{\R}{\text{R}}
\newcommand{\M}{\text{M}}
\newcommand{\m}{\text{m}}
\newcommand{\s}{\text{s}}
\newcommand{\f}{\text{f}}
\newcommand{\K}{\text{K}}
\newcommand{\hc}{\text{h.c.}}
\newcommand{\kf}{k_\F}
\newcommand{\vf}{v_\F}
\newcommand{\const}{\text{const.}}
\renewcommand{\L}{\text{L}}
\newcommand{\gray}[1]{\bgroup\color{white!40!black}#1\egroup}
\newcommand{\TODO}[1]{\bgroup\color{red!60!black}#1\egroup}
\newcommand\QnA[1]{\bgroup\color{blue}#1\egroup}
\newcommand\highlight[1]{\bgroup\color{green!60!black}#1\egroup}

\newcommand{\cabx}[3]{\begin{smallmatrix}\gray{A}&\gray{:}&#1 \\ \gray{B}&\gray{:}&#2 \\ \gray{\Xi}&\gray{:}&#3\end{smallmatrix}}
\newcommand{\zabx}[3]{\begin{smallmatrix}\gray{A}&\gray{:}&#1 \\ \gray{B}&\gray{:}&#2 \\ \gray{Z}&\gray{:}&#3\end{smallmatrix}}
\newcommand{\abx}[3]{\begin{smallmatrix}#1 \\ #2 \\ #3\end{smallmatrix}}

\newcommand{\normord}[1]{
  {:\mathrel{\mspace{1mu}#1\mspace{1mu}}:}
}
\newcommand{\varnormord}[1]{
  {\mathrel{\mspace{1mu}\substack{\ast\\[-1.5pt] \ast}#1\substack{\ast\\[-1.5pt] \ast}\mspace{1mu}}}
}

% punctuation styles (see https://tex.stackexchange.com/questions/56063/how-to-properly-typeset-footnotes-superscripts-after-punctuation-marks)
\let\origfootnote\footnote
\renewcommand{\footnote}[1]{\kern.06em\origfootnote{#1}}
\newcommand{\punctfootnote}[1]{\kern-.06em\origfootnote{#1}}

% Multiline equations
\newcommand{\bsplit}[1]{\begin{aligned}[t]#1\end{aligned}}